\section{The Foundation}

A reasonable concern about any candidate theory of everything is that it carries too many free dials. If the foundation could have been assembled in a hundred ways, selecting one looks arbitrary. We therefore ask, of the five base elements, how many are genuine choices and how many are forced by the demands placed on them.

The result is that four of the five are heavily constrained, several of them to a single option, and only one, the arrow, is a deliberate posit. What follows is a filtering argument, not a proof that the universe must take this form. We mark each step as a theorem or as a preference.

The five base elements are the following.

\begin{itemize}
  \item the mesh ($\{5,3,4\}$)
  \item the ternary tone ($-1, 0, +1$)
  \item the conserving perception rule
  \item eternal growth by reflection
  \item the arrow (the value-direction)
\end{itemize}

Their full definitions appear in the next section.

\subsection{The mesh narrows to the dodecagrid}

Three considerations reduce the space of meshes to one, with the final step a stated preference rather than a theorem.

Curvature comes first. A space of uniform curvature is spherical, flat, or hyperbolic. A spherical space closes up and stops growing. A flat space offers only polynomial room, and as a regular lattice it carries a preferred frame that breaks Lorentz invariance. A hyperbolic space is the one option that is infinite, grows exponentially, and carries no preferred direction. We take hyperbolic curvature as forced by the demands of unbounded growth and frame independence.

Dimension comes next. Compact regular hyperbolic honeycombs exist only in dimensions two, three, and four. This is a theorem, established by experiment (P62). Among those, only three dimensions support stable closed orbits, so only in three dimensions do bound configurations hold together (P68). That leaves dimension three.

The choice of crystal comes last. Three dimensions admit four regular hyperbolic honeycombs. Ranking them by ternary cell vertices, by the golden pentagonal cell, and by the tightest eternal closure points to $\{5,3,4\}$, the dodecagrid. This step uses a minimality preference, not a theorem. The auto-selection experiment (P86) confirms that the ternary tone together with minimal eternal closure lands on $\{5,3,4\}$.

\begin{figure}[h]
\centering
\includegraphics[width=\textwidth]{vibe-mesh-534.png}
\caption{The chosen crystal, the dodecagrid $\{5,3,4\}$, dodecahedral cells with four meeting around every edge, filling curved three-dimensional space. This is the substrate the rest of the paper builds on.}
\label{fig:dodecagrid-chosen}
\end{figure}

The mesh is therefore forced to be hyperbolic and three-dimensional by real constraints, and to $\{5,3,4\}$ specifically by a minimality preference among the few remaining options. We present this as strong, not airtight, and we say so plainly.

\subsection{The tone narrows to ternary}

How many felt states should a vibe carry? One state makes no distinction. Two states without a neutral give opposition but leave the rule's ties nowhere to rest. Four or more states are redundant for a sign. The minimal signed alphabet that supports opposition together with a rest state is three, namely $\{-1, 0, +1\}$. This also matches the natural output of a sign function, which returns negative, zero, or positive. The argument is a minimality argument, and it is a clean one.

\subsection{The rule is tightly constrained}

The update faces three demands.

\begin{itemize}
  \item It must be local. A vibe reads only its neighbors, with no global clock.
  \item It must treat the two charges symmetrically, so that flipping every tone flips the rule with it.
  \item In the committed model it must conserve the total charge, so that creation is balanced.
\end{itemize}

The conserving perception rule, in which opposite tones annihilate to peace, a charge hops into an empty neighbor, and peace spawns a balanced pair under the arrow, is the simplest local, charge-symmetric, conserving rule of this kind. It has no free parameters beyond the move rates, and it has no hidden state. The only state is the tone. The note is perception, what a vibe reads of its neighbors, not a stored quantity.

An earlier version of the work used a non-conserving rule based on the sign of a neighbor vote. The move to the conserving exchange was made to support a genuine conserved charge and the reversibility that later proved to be the precondition for a real quantum field. This element was refined as the work progressed, and the current rule is the one the experiments use.

\subsection{Growth is reflection}

Space must come from somewhere and keep coming. The minimal way to repeat a distinction is to mirror it. A reflection performed twice returns to the start, which makes it the smallest unit of repetition. Reflections are the building blocks of all rigid motion, by the Cartan-Dieudonn\'e theorem. Reflecting consistently, with neither gaps nor overlaps, forces the mirror angles to be whole-number fractions, which is precisely a Coxeter group, the crystal. Eternal growth by reflection is therefore not an extra mechanism bolted on. It is the same machinery that builds the tiling.

\subsection{From nothing to something}

There was most likely never a beginning. A true nothing is inherently unstable, an impossibility rather than a starting point. A perfect absence has nothing in it to enforce its own emptiness, so the smallest distinction is enough to break it, and a distinction is always available. The right question is not how something arose from nothing at a first moment. It is why there is a stable something rather than an unstable nothing, and the answer is that nothing cannot persist while a self-consistent something can.

Given a single forced distinction, the reflection machinery above does the rest. The distinction is mirrored, the mirrorings compose into a Coxeter group, and the crystal grows outward without end. No first cell is required. The growth is eternal both forward in time and as an explanation, since the crystal is its own ongoing reason to continue. Table~\ref{tab:bootstrap} steps through the bootstrap.

\begin{table}[h]
\centering
\begin{tabular}{cll}
\toprule
Step & Configuration & What forces the next step \\
\midrule
0 & Nothing & Unstable. A perfect absence cannot enforce itself. \\
1 & One distinction & The minimal asymmetry, the seed of structure. \\
2 & A mirror of it & The simplest way to repeat a distinction. \\
3 & Consistent reflection & Whole-number mirror angles, a Coxeter group, the crystal. \\
4 and on & Eternal growth & New cells born at the frontier without end. \\
\bottomrule
\end{tabular}
\caption{The bootstrap from nothing to a growing crystal. Nothing is not a stable initial state but an impossibility, so the sequence has no first moment that needs explaining, only a reason that something self-consistent persists.}
\label{tab:bootstrap}
\end{table}

\subsection{The arrow is the one deliberate posit}

The fifth element differs in kind. The arrow, the value-direction in which pain is less than peace which is less than pleasure, with desire leaning toward pleasure, is not forced by any minimality argument. It is a value, and values are chosen rather than derived. We add it by decision, and we make it conserving, so that it creates only balanced pairs and coexists with charge conservation.

Without it the universe relaxes to dead peace. With it the universe stays alive (P100). The chapters on life, mind, and intelligence trace back to the arrow together with the four constrained elements. The accounting is therefore four base elements heavily constrained and one base element deliberately posited.

\subsection{What this does and does not claim}

The claim is that the foundation has very few real dials, and that the structural ones, hyperbolic curvature, three dimensions, the ternary tone, conservation, and reflection, follow from modest demands rather than from taste. This matters, because it means the theory does not buy its breadth with a long list of tuned choices. The fine-tuning test (P166) supports this at the level of the rule's rates. The rich behavior fills most of the parameter space rather than a narrow window.

The claim is not that the foundation is logically inevitable. The choice of crystal is a minimality preference, the arrow is a chosen value, and the whole argument is a filter, not a derivation from nothing. We present it as a strong reduction of arbitrariness, and no more.
